\documentclass[10pt]{article}
% Hojas resumen: un A4 por delante y por detrás, dos columnas por cara.
% Uso: \documentclass[10pt]{article} % Hojas resumen: un A4 por delante y por detrás, dos columnas por cara.
% Uso: \documentclass[10pt]{article} % Hojas resumen: un A4 por delante y por detrás, dos columnas por cara.
% Uso: \documentclass[10pt]{article} \input{estilo-resumen.tex} \bloque{analisis|algebra|azar} \cabecera{n}{Título}{Lecciones}
\usepackage[utf8]{inputenc}
\usepackage[T1]{fontenc}
\usepackage[spanish,es-nodecimaldot]{babel}
\usepackage{lmodern}
\usepackage{amsmath,amssymb}
\usepackage{graphicx}
\usepackage[a4paper,left=10mm,right=10mm,top=15mm,bottom=13mm,headheight=11pt,headsep=3mm,footskip=6mm]{geometry}
\usepackage{multicol}
\setlength{\columnsep}{6mm}
\setlength{\columnseprule}{0.3pt}
\usepackage{xcolor}
\usepackage{enumitem}
\setlist{nosep,leftmargin=1.1em,itemsep=1pt}
\usepackage{array}
\usepackage{needspace}
\usepackage{pgfplots}
\usepgfplotslibrary{fillbetween}
\pgfplotsset{compat=1.18,
  mini/.style={width=4.3cm,height=3.3cm,scale only axis=false,axis lines=middle,xtick=\empty,ytick=\empty,samples=120,
    xlabel={\scriptsize$x$},ylabel={\scriptsize$y$},title style={font=\scriptsize,yshift=-4pt},clip=true,
    every axis plot/.append style={thick,color=acento}}}
\usepackage{ragged2e}
\usepackage[most]{tcolorbox}
\usepackage{fancyhdr,lastpage}
\usepackage[hidelinks,pdfencoding=auto]{hyperref}
\hypersetup{pdflang=es-ES,pdfsubject={Matemáticas Aplicadas a las Ciencias Sociales II}}

\definecolor{analisis}{HTML}{1D4ED8}
\definecolor{algebra}{HTML}{86198F}
\definecolor{azar}{HTML}{9A3412}
\definecolor{aviso}{HTML}{B91C1C}
\colorlet{acento}{analisis}
\newcommand{\bloque}[1]{\colorlet{acento}{#1}}

\setlength{\parindent}{0pt}
\setlength{\parskip}{2pt}
\renewcommand{\arraystretch}{1.2}
\AtBeginDocument{\RaggedRight\fontsize{9.8}{11.9}\selectfont\setlength{\abovedisplayskip}{3pt}\setlength{\belowdisplayskip}{3pt}\setlength{\abovedisplayshortskip}{1pt}\setlength{\belowdisplayshortskip}{1pt}}

\newcommand{\temanum}{}\newcommand{\tematitulo}{}
\pagestyle{fancy}
\fancyhf{}
\fancyhead[L]{\footnotesize\color{acento}\textbf{Hoja resumen}\quad Tema \temanum. \tematitulo}
\fancyhead[R]{\footnotesize Matemáticas Aplicadas a las Ciencias Sociales II --- 2.\textsuperscript{o} Bachillerato}
\fancyfoot[L]{\scriptsize IES Al-Ándalus --- Juan María Gámez Ortiz, con la ayuda de Claude y de Codex --- CC BY-NC 4.0}
\fancyfoot[R]{\scriptsize Cara \thepage\ de \pageref{LastPage}}
\renewcommand{\headrulewidth}{0.4pt}

% Cabecera de la primera cara: número, título y lecciones del tema.
\newcommand{\cabecera}[3]{%
  \renewcommand{\temanum}{#1}\renewcommand{\tematitulo}{#2}%
  \hypersetup{pdftitle={Hoja resumen. Tema #1. #2}}%
  \noindent\begin{tcolorbox}[enhanced,colback=acento!8,colframe=acento,boxrule=0.6pt,arc=2pt,left=4pt,right=4pt,top=2pt,bottom=2pt]
  \begin{minipage}[c]{1.1cm}\includegraphics[width=1cm]{logo-al-andalus.png}\end{minipage}%
  \begin{minipage}[c]{\dimexpr\linewidth-1.2cm}{\large\bfseries\color{acento}Tema #1. #2}\\[1pt]{\scriptsize #3}\end{minipage}
  \end{tcolorbox}\vspace{-2pt}}

% Apartado con filete del color del bloque.
\newcommand{\seccion}[1]{\par\needspace{5\baselineskip}\vspace{3pt}{\color{acento}\bfseries\normalsize #1}\par\vspace{-1pt}{\color{acento}\hrule height 0.5pt}\vspace{2pt}}
\newcommand{\sub}[1]{\par\needspace{3\baselineskip}\vspace{1pt}\textbf{\color{acento}#1}\enspace}

% Fórmula clave, aviso de errores y ejemplo.
\tcbset{before upper=\RaggedRight}
\newtcolorbox{formula}{enhanced,colback=acento!6,colframe=acento!70,boxrule=0.5pt,arc=1.5pt,left=3pt,right=3pt,top=0pt,bottom=0pt,before skip=2pt,after skip=2pt}
\newtcolorbox{ojo}{enhanced,colback=aviso!5,colframe=aviso!60,boxrule=0.5pt,arc=1.5pt,left=3pt,right=3pt,top=1pt,bottom=1pt,before skip=3pt,after skip=2pt,
  fonttitle=\bfseries\footnotesize,coltitle=aviso,colbacktitle=aviso!5,title={Ojo},attach title to upper={\enspace}}
\newtcolorbox{ejemplo}{enhanced,colback=white,colframe=acento!45,boxrule=0.4pt,arc=1.5pt,left=3pt,right=3pt,top=1pt,bottom=1pt,before skip=3pt,after skip=2pt,
  fonttitle=\bfseries\footnotesize,coltitle=acento,colbacktitle=white,title={Ejemplo},attach title to upper={\enspace}}
\newcommand{\dx}{\,dx}
 \bloque{analisis|algebra|azar} \cabecera{n}{Título}{Lecciones}
\usepackage[utf8]{inputenc}
\usepackage[T1]{fontenc}
\usepackage[spanish,es-nodecimaldot]{babel}
\usepackage{lmodern}
\usepackage{amsmath,amssymb}
\usepackage{graphicx}
\usepackage[a4paper,left=10mm,right=10mm,top=15mm,bottom=13mm,headheight=11pt,headsep=3mm,footskip=6mm]{geometry}
\usepackage{multicol}
\setlength{\columnsep}{6mm}
\setlength{\columnseprule}{0.3pt}
\usepackage{xcolor}
\usepackage{enumitem}
\setlist{nosep,leftmargin=1.1em,itemsep=1pt}
\usepackage{array}
\usepackage{needspace}
\usepackage{pgfplots}
\usepgfplotslibrary{fillbetween}
\pgfplotsset{compat=1.18,
  mini/.style={width=4.3cm,height=3.3cm,scale only axis=false,axis lines=middle,xtick=\empty,ytick=\empty,samples=120,
    xlabel={\scriptsize$x$},ylabel={\scriptsize$y$},title style={font=\scriptsize,yshift=-4pt},clip=true,
    every axis plot/.append style={thick,color=acento}}}
\usepackage{ragged2e}
\usepackage[most]{tcolorbox}
\usepackage{fancyhdr,lastpage}
\usepackage[hidelinks,pdfencoding=auto]{hyperref}
\hypersetup{pdflang=es-ES,pdfsubject={Matemáticas Aplicadas a las Ciencias Sociales II}}

\definecolor{analisis}{HTML}{1D4ED8}
\definecolor{algebra}{HTML}{86198F}
\definecolor{azar}{HTML}{9A3412}
\definecolor{aviso}{HTML}{B91C1C}
\colorlet{acento}{analisis}
\newcommand{\bloque}[1]{\colorlet{acento}{#1}}

\setlength{\parindent}{0pt}
\setlength{\parskip}{2pt}
\renewcommand{\arraystretch}{1.2}
\AtBeginDocument{\RaggedRight\fontsize{9.8}{11.9}\selectfont\setlength{\abovedisplayskip}{3pt}\setlength{\belowdisplayskip}{3pt}\setlength{\abovedisplayshortskip}{1pt}\setlength{\belowdisplayshortskip}{1pt}}

\newcommand{\temanum}{}\newcommand{\tematitulo}{}
\pagestyle{fancy}
\fancyhf{}
\fancyhead[L]{\footnotesize\color{acento}\textbf{Hoja resumen}\quad Tema \temanum. \tematitulo}
\fancyhead[R]{\footnotesize Matemáticas Aplicadas a las Ciencias Sociales II --- 2.\textsuperscript{o} Bachillerato}
\fancyfoot[L]{\scriptsize IES Al-Ándalus --- Juan María Gámez Ortiz, con la ayuda de Claude y de Codex --- CC BY-NC 4.0}
\fancyfoot[R]{\scriptsize Cara \thepage\ de \pageref{LastPage}}
\renewcommand{\headrulewidth}{0.4pt}

% Cabecera de la primera cara: número, título y lecciones del tema.
\newcommand{\cabecera}[3]{%
  \renewcommand{\temanum}{#1}\renewcommand{\tematitulo}{#2}%
  \hypersetup{pdftitle={Hoja resumen. Tema #1. #2}}%
  \noindent\begin{tcolorbox}[enhanced,colback=acento!8,colframe=acento,boxrule=0.6pt,arc=2pt,left=4pt,right=4pt,top=2pt,bottom=2pt]
  \begin{minipage}[c]{1.1cm}\includegraphics[width=1cm]{logo-al-andalus.png}\end{minipage}%
  \begin{minipage}[c]{\dimexpr\linewidth-1.2cm}{\large\bfseries\color{acento}Tema #1. #2}\\[1pt]{\scriptsize #3}\end{minipage}
  \end{tcolorbox}\vspace{-2pt}}

% Apartado con filete del color del bloque.
\newcommand{\seccion}[1]{\par\needspace{5\baselineskip}\vspace{3pt}{\color{acento}\bfseries\normalsize #1}\par\vspace{-1pt}{\color{acento}\hrule height 0.5pt}\vspace{2pt}}
\newcommand{\sub}[1]{\par\needspace{3\baselineskip}\vspace{1pt}\textbf{\color{acento}#1}\enspace}

% Fórmula clave, aviso de errores y ejemplo.
\tcbset{before upper=\RaggedRight}
\newtcolorbox{formula}{enhanced,colback=acento!6,colframe=acento!70,boxrule=0.5pt,arc=1.5pt,left=3pt,right=3pt,top=0pt,bottom=0pt,before skip=2pt,after skip=2pt}
\newtcolorbox{ojo}{enhanced,colback=aviso!5,colframe=aviso!60,boxrule=0.5pt,arc=1.5pt,left=3pt,right=3pt,top=1pt,bottom=1pt,before skip=3pt,after skip=2pt,
  fonttitle=\bfseries\footnotesize,coltitle=aviso,colbacktitle=aviso!5,title={Ojo},attach title to upper={\enspace}}
\newtcolorbox{ejemplo}{enhanced,colback=white,colframe=acento!45,boxrule=0.4pt,arc=1.5pt,left=3pt,right=3pt,top=1pt,bottom=1pt,before skip=3pt,after skip=2pt,
  fonttitle=\bfseries\footnotesize,coltitle=acento,colbacktitle=white,title={Ejemplo},attach title to upper={\enspace}}
\newcommand{\dx}{\,dx}
 \bloque{analisis|algebra|azar} \cabecera{n}{Título}{Lecciones}
\usepackage[utf8]{inputenc}
\usepackage[T1]{fontenc}
\usepackage[spanish,es-nodecimaldot]{babel}
\usepackage{lmodern}
\usepackage{amsmath,amssymb}
\usepackage{graphicx}
\usepackage[a4paper,left=10mm,right=10mm,top=15mm,bottom=13mm,headheight=11pt,headsep=3mm,footskip=6mm]{geometry}
\usepackage{multicol}
\setlength{\columnsep}{6mm}
\setlength{\columnseprule}{0.3pt}
\usepackage{xcolor}
\usepackage{enumitem}
\setlist{nosep,leftmargin=1.1em,itemsep=1pt}
\usepackage{array}
\usepackage{needspace}
\usepackage{pgfplots}
\usepgfplotslibrary{fillbetween}
\pgfplotsset{compat=1.18,
  mini/.style={width=4.3cm,height=3.3cm,scale only axis=false,axis lines=middle,xtick=\empty,ytick=\empty,samples=120,
    xlabel={\scriptsize$x$},ylabel={\scriptsize$y$},title style={font=\scriptsize,yshift=-4pt},clip=true,
    every axis plot/.append style={thick,color=acento}}}
\usepackage{ragged2e}
\usepackage[most]{tcolorbox}
\usepackage{fancyhdr,lastpage}
\usepackage[hidelinks,pdfencoding=auto]{hyperref}
\hypersetup{pdflang=es-ES,pdfsubject={Matemáticas Aplicadas a las Ciencias Sociales II}}

\definecolor{analisis}{HTML}{1D4ED8}
\definecolor{algebra}{HTML}{86198F}
\definecolor{azar}{HTML}{9A3412}
\definecolor{aviso}{HTML}{B91C1C}
\colorlet{acento}{analisis}
\newcommand{\bloque}[1]{\colorlet{acento}{#1}}

\setlength{\parindent}{0pt}
\setlength{\parskip}{2pt}
\renewcommand{\arraystretch}{1.2}
\AtBeginDocument{\RaggedRight\fontsize{9.8}{11.9}\selectfont\setlength{\abovedisplayskip}{3pt}\setlength{\belowdisplayskip}{3pt}\setlength{\abovedisplayshortskip}{1pt}\setlength{\belowdisplayshortskip}{1pt}}

\newcommand{\temanum}{}\newcommand{\tematitulo}{}
\pagestyle{fancy}
\fancyhf{}
\fancyhead[L]{\footnotesize\color{acento}\textbf{Hoja resumen}\quad Tema \temanum. \tematitulo}
\fancyhead[R]{\footnotesize Matemáticas Aplicadas a las Ciencias Sociales II --- 2.\textsuperscript{o} Bachillerato}
\fancyfoot[L]{\scriptsize IES Al-Ándalus --- Juan María Gámez Ortiz, con la ayuda de Claude y de Codex --- CC BY-NC 4.0}
\fancyfoot[R]{\scriptsize Cara \thepage\ de \pageref{LastPage}}
\renewcommand{\headrulewidth}{0.4pt}

% Cabecera de la primera cara: número, título y lecciones del tema.
\newcommand{\cabecera}[3]{%
  \renewcommand{\temanum}{#1}\renewcommand{\tematitulo}{#2}%
  \hypersetup{pdftitle={Hoja resumen. Tema #1. #2}}%
  \noindent\begin{tcolorbox}[enhanced,colback=acento!8,colframe=acento,boxrule=0.6pt,arc=2pt,left=4pt,right=4pt,top=2pt,bottom=2pt]
  \begin{minipage}[c]{1.1cm}\includegraphics[width=1cm]{logo-al-andalus.png}\end{minipage}%
  \begin{minipage}[c]{\dimexpr\linewidth-1.2cm}{\large\bfseries\color{acento}Tema #1. #2}\\[1pt]{\scriptsize #3}\end{minipage}
  \end{tcolorbox}\vspace{-2pt}}

% Apartado con filete del color del bloque.
\newcommand{\seccion}[1]{\par\needspace{5\baselineskip}\vspace{3pt}{\color{acento}\bfseries\normalsize #1}\par\vspace{-1pt}{\color{acento}\hrule height 0.5pt}\vspace{2pt}}
\newcommand{\sub}[1]{\par\needspace{3\baselineskip}\vspace{1pt}\textbf{\color{acento}#1}\enspace}

% Fórmula clave, aviso de errores y ejemplo.
\tcbset{before upper=\RaggedRight}
\newtcolorbox{formula}{enhanced,colback=acento!6,colframe=acento!70,boxrule=0.5pt,arc=1.5pt,left=3pt,right=3pt,top=0pt,bottom=0pt,before skip=2pt,after skip=2pt}
\newtcolorbox{ojo}{enhanced,colback=aviso!5,colframe=aviso!60,boxrule=0.5pt,arc=1.5pt,left=3pt,right=3pt,top=1pt,bottom=1pt,before skip=3pt,after skip=2pt,
  fonttitle=\bfseries\footnotesize,coltitle=aviso,colbacktitle=aviso!5,title={Ojo},attach title to upper={\enspace}}
\newtcolorbox{ejemplo}{enhanced,colback=white,colframe=acento!45,boxrule=0.4pt,arc=1.5pt,left=3pt,right=3pt,top=1pt,bottom=1pt,before skip=3pt,after skip=2pt,
  fonttitle=\bfseries\footnotesize,coltitle=acento,colbacktitle=white,title={Ejemplo},attach title to upper={\enspace}}
\newcommand{\dx}{\,dx}

\bloque{analisis}
\begin{document}
\cabecera{1}{Límites y derivadas}{Lecciones 1.1--1.18: límites, operaciones con el infinito, indeterminaciones, continuidad, derivada, rectas tangente y normal, tabla y reglas de derivación, regla de L'Hôpital.}
\begin{multicols}{2}

\seccion{1. Límite de una función}
\(\lim_{x\to a}f(x)=L\) significa que \(f(x)\) se acerca a \(L\) cuando \(x\) se acerca a \(a\), sin tomar \(x=a\). El valor \(f(a)\) no interviene: puede no existir o ser distinto de \(L\).
\begin{formula}
\[\lim_{x\to a}f(x)=L\iff \underbrace{\lim_{x\to a^-}f(x)}_{L_-}=\underbrace{\lim_{x\to a^+}f(x)}_{L_+}=L.\]
\end{formula}
\sub{Laterales.} Imprescindibles en funciones a trozos, en valores absolutos y cuando un denominador se anula: el signo puede ser distinto a cada lado.
\begin{ejemplo}
\[f(x)=\begin{cases}x+1,&x<2\\ x^2-1,&x\ge2\end{cases}\]
\(\lim_{x\to2^-}f=3\) y \(\lim_{x\to2^+}f=3\), así que \(\lim_{x\to2}f(x)=3\).
\end{ejemplo}

\seccion{2. Operaciones con el infinito}
\begin{tabular}{@{}ll@{}}
\(\infty+\infty=\infty\), \(\;k\pm\infty=\pm\infty\) & \(\infty\cdot\infty=\infty\)\\
\(k\cdot\infty=\pm\infty\) \((k\ne0)\) & \(\dfrac{k}{\pm\infty}=0\)\\
\(\dfrac{k}{0}=\pm\infty\) \((k\ne0\); signo por laterales\()\) & \(\dfrac{\infty}{k}=\pm\infty\)\\
\(a^{+\infty}=+\infty\) si \(a>1\); \(\;=0\) si \(0\le a<1\) & \(\infty^{k}=\infty\) \((k>0)\)
\end{tabular}
\begin{formula}
\textbf{Indeterminaciones:}
\[\frac00,\quad \frac\infty\infty,\quad \infty-\infty,\quad 0\cdot\infty,\quad 1^\infty,\quad 0^0,\quad \infty^0\]
\end{formula}
Una indeterminación no es un resultado: hay que transformar la expresión.

\seccion{3. Cálculo de límites}
\begin{enumerate}
  \item \textbf{Sustituir.} Si sale un número, ese es el límite.
  \item \(\dfrac k0\): estudiar el signo de cada lateral \(\Rightarrow\pm\infty\).
  \item \(\dfrac00\) con polinomios: \textbf{factorizar} y simplificar \((x-a)\).
  \item \(\dfrac00\) con raíces: multiplicar por el \textbf{conjugado}.
  \item \(\dfrac\infty\infty\) con polinomios (\(x\to\pm\infty\)): comparar \textbf{grados}.
  \item \(\infty-\infty\): común denominador o conjugado.
  \item \(1^\infty\): \(\lim f^g=e^{\lim g\,(f-1)}\).
\end{enumerate}
\begin{formula}
\(\displaystyle\lim_{x\to\pm\infty}\frac{a_nx^n+\cdots}{b_mx^m+\cdots}=\begin{cases}\pm\infty & n>m\ (\text{signo}),\\[1pt] \dfrac{a_n}{b_m} & n=m,\\[4pt] 0 & n<m.\end{cases}\)
\end{formula}
\begin{ejemplo}
\[\lim_{x\to2}\frac{x^2-4}{x-2}=\Bigl[\frac00\Bigr]=\lim_{x\to2}\frac{(x-2)(x+2)}{x-2}=4.\]
\[\lim_{x\to+\infty}\bigl(\sqrt{x^2+x}-x\bigr)=\lim_{x\to+\infty}\frac{x}{\sqrt{x^2+x}+x}=\frac12.\]
\[\lim_{x\to+\infty}\Bigl(1+\frac2x\Bigr)^{x}=[1^\infty]=e^{\lim x\cdot\frac2x}=e^2.\]
\end{ejemplo}
\sub{Crecimiento en \(+\infty\).} \(\ln x\ll x^n\ll e^x\ (n>0)\): por ejemplo, \(\lim\frac{x^3}{e^x}=0\) y \(\lim\frac{\ln x}{x}=0\).

\seccion{4. Asíntotas}
\begin{itemize}
  \item \textbf{Vertical} \(x=a\): algún lateral \(\lim_{x\to a^\pm}f(x)=\pm\infty\). Candidatos: ceros del denominador, borde del dominio del logaritmo.
  \item \textbf{Horizontal} \(y=L\): \(\lim_{x\to\pm\infty}f(x)=L\) (finito). Se estudia en \(+\infty\) y en \(-\infty\).
  \item \textbf{Oblicua} \(y=mx+n\): \(m=\lim\dfrac{f(x)}x\ne0\) finito, \(n=\lim\bigl(f(x)-mx\bigr)\). No coexiste con una horizontal en el mismo lado.
\end{itemize}
\begin{ejemplo}
\(f(x)=\dfrac{2x+1}{x-3}\). Como \(\lim_{x\to3^\pm}f=\pm\infty\), hay asíntota vertical \(x=3\). Como \(\lim_{x\to\pm\infty}f=2\), hay asíntota horizontal \(y=2\).
\end{ejemplo}

\seccion{5. Continuidad}
\begin{formula}
\(f\) es continua en \(a\iff\) existe \(f(a)\), existe \(\lim_{x\to a}f(x)\) y \(\lim_{x\to a^-}f=\lim_{x\to a^+}f=f(a).\)
\end{formula}
\textbf{Discontinuidades:} \emph{evitable} (el límite existe pero no coincide con \(f(a)\) o no existe \(f(a)\)); \emph{de salto finito} (laterales finitos y distintos); \emph{de salto infinito} (algún lateral infinito).
Polinomios, exponenciales, \(\sin\), \(\cos\) son continuas en \(\mathbb R\); racionales y logaritmos, en su dominio. En funciones a trozos se estudian los puntos de unión.
\begin{ejemplo}
\[f(x)=\begin{cases}x^2+k,&x<1\\ 3x,&x\ge1\end{cases}\]
es continua en 1 si \(1+k=3\), es decir, \(k=2\).
\end{ejemplo}

\seccion{6. La derivada}
\begin{formula}
\[\operatorname{TVM}_{[a,a+h]}=\frac{f(a+h)-f(a)}h\]
\[f'(a)=\lim_{h\to0}\frac{f(a+h)-f(a)}h\]
\end{formula}
\(f'(a)\) es la \textbf{pendiente de la recta tangente} y la \textbf{tasa de variación instantánea}. \(f\) es derivable en \(a\) si \(f'(a^-)=f'(a^+)\) (derivadas laterales finitas e iguales).
\begin{ojo}
Derivable \(\Rightarrow\) continua, pero no al revés: \(|x|\) es continua en 0 y no derivable (pico: \(f'(0^-)=-1\ne1=f'(0^+)\)). Si no es continua, no es derivable.
\end{ojo}

\seccion{7. Rectas tangente y normal en \(x=a\)}
\begin{formula}
\[\text{Tangente: }\ y-f(a)=f'(a)\,(x-a)\]
\[\text{Normal: }\ y-f(a)=-\frac1{f'(a)}\,(x-a)\]
\end{formula}
Pasos: \(f(a)\), \(f'(x)\), \(f'(a)\) y sustituir. Si \(f'(a)=0\), la tangente es horizontal \(y=f(a)\) y la normal es vertical \(x=a\).
\begin{ejemplo}
\(f(x)=x^2-3x\) en \(a=1\): \(f(1)=-2\), \(f'(x)=2x-3\), \(f'(1)=-1\).
Tangente: \(y+2=-(x-1)\), es decir, \(y=-x-1\). Normal: \(y+2=x-1\), es decir, \(y=x-3\).
\end{ejemplo}

\seccion{8. Tabla de derivadas}
\begin{tabular}{@{}l@{\ \ }l@{}}
\textbf{Función} & \textbf{Compuesta} (\(u=u(x)\))\\ \hline
\((k)'=0\) & \\
\((x^n)'=nx^{n-1}\) & \((u^n)'=n\,u^{n-1}u'\)\\
\((\sqrt x)'=\dfrac1{2\sqrt x}\) & \((\sqrt u)'=\dfrac{u'}{2\sqrt u}\)\\
\((e^x)'=e^x\) & \((e^u)'=u'\,e^u\)\\
\((a^x)'=a^x\ln a\) & \((a^u)'=u'\,a^u\ln a\)\\
\((\ln x)'=\dfrac1x\) & \((\ln u)'=\dfrac{u'}u\)\\
\((\log_a x)'=\dfrac1{x\ln a}\) & \((\log_a u)'=\dfrac{u'}{u\ln a}\)\\
\((\sin x)'=\cos x\) & \((\sin u)'=u'\cos u\)\\
\((\cos x)'=-\sin x\) & \((\cos u)'=-u'\sin u\)\\
\((\operatorname{tg}x)'=\dfrac1{\cos^2x}\) & \((\operatorname{tg}u)'=\dfrac{u'}{\cos^2u}\)
\end{tabular}

\seccion{9. Reglas de derivación}
\begin{formula}
\((f\pm g)'=f'\pm g'\qquad (k f)'=k f'\qquad (fg)'=f'g+fg'\)\\[2pt]
\(\Bigl(\dfrac fg\Bigr)'=\dfrac{f'g-fg'}{g^2}\ (g\ne0)\qquad \bigl(F(g(x))\bigr)'=F'(g(x))\,g'(x)\)
\end{formula}
\textbf{Regla de la cadena:} se deriva la función exterior evaluada en la interior y se multiplica por la derivada de la interior.
\begin{ejemplo}
\[\Bigl(\frac{x^2}{x+1}\Bigr)'=\frac{2x(x+1)-x^2}{(x+1)^2}=\frac{x^2+2x}{(x+1)^2}\]
\[(x\ln x)'=\ln x+1,\qquad \bigl(\ln(x^2+1)\bigr)'=\frac{2x}{x^2+1}\]
\end{ejemplo}
\begin{ejemplo}
\[\bigl((x^2+1)^3\bigr)'=3(x^2+1)^2\cdot2x=6x(x^2+1)^2\]
\[\bigl(e^{-x^2}\bigr)'=-2x\,e^{-x^2}\]
\end{ejemplo}

\sub{Derivabilidad de una función a trozos.} Primero continuidad en el punto de unión; después, derivadas laterales iguales.
\begin{ejemplo}
\[f(x)=\begin{cases}x^2,&x\le1\\ 2x-1,&x>1\end{cases}\]
Es continua (los dos trozos valen 1) y las derivadas laterales coinciden, \(f'(1^-)=f'(1^+)=2\): derivable en 1.
\end{ejemplo}

\seccion{10. Regla de L'Hôpital}
\begin{formula}
Si \(\lim\frac{f}{g}\) da \(\frac00\) o \(\frac\infty\infty\) y existe \(\lim\frac{f'}{g'}\):
\[\lim\frac{f(x)}{g(x)}=\lim\frac{f'(x)}{g'(x)}\]
\end{formula}
Se derivan \textbf{numerador y denominador por separado} (no es la derivada del cociente). Se puede repetir si sigue indeterminado. \(0\cdot\infty\) e \(\infty-\infty\) se transforman antes en cociente.
\begin{ejemplo}
\[\lim_{x\to0}\frac{e^x-1}{x}=\Bigl[\frac00\Bigr]=\lim_{x\to0}\frac{e^x}{1}=1\]
\[\lim_{x\to0}\frac{1-\cos x}{x^2}=\lim_{x\to0}\frac{\sin x}{2x}=\lim_{x\to0}\frac{\cos x}{2}=\frac12\]
\[\lim_{x\to+\infty}\frac{x^2}{e^x}=\lim\frac{2x}{e^x}=\lim\frac{2}{e^x}=0\]
\end{ejemplo}

\seccion{11. La derivada en economía}
La derivada mide la \textbf{variación instantánea}. Si \(C(q)\) es el coste de producir \(q\) unidades, \(C'(q)\) es el \textbf{coste marginal}: aproxima el coste de producir una unidad más. Igual con ingreso \(I'(q)\) y beneficio \(B'(q)=I'(q)-C'(q)\).
\begin{ejemplo}
\(C(q)=q^2+10q+50\) euros: \(C'(q)=2q+10\), \(C'(20)=50\). Con 20 unidades, la siguiente cuesta unos 50\,euros (exactamente, \(C(21)-C(20)=51\)).
\end{ejemplo}

\seccion{12. Ejercicio guiado: parámetros}
Halla \(a\) y \(b\) para que sea derivable en \(x=1\):
\[f(x)=\begin{cases}ax^2+1,&x\le1\\ bx+3,&x>1\end{cases}\]
\begin{enumerate}
  \item \textbf{Continuidad:} \(\lim_{x\to1^-}f=a+1\), \(\lim_{x\to1^+}f=b+3\Rightarrow a+1=b+3\).
  \item \textbf{Derivadas laterales:} \(f'(1^-)=2a\), \(f'(1^+)=b\Rightarrow 2a=b\).
  \item \textbf{Resolver:} \(a+1=2a+3\Rightarrow a=-2\), \(b=-4\).
  \item \textbf{Comprobar:} \(f(1)=-1\) por ambos lados y \(f'(1)=-4\) por ambos lados.
\end{enumerate}

\seccion{13. Estrategia ante un límite}
\begin{enumerate}
  \item Sustituye y clasifica el resultado: número, \(\frac k0\) o indeterminación.
  \item Elige la herramienta según el tipo (apartado 3) o L'Hôpital.
  \item Si \(x\to a\) con cambio de signo o trozos, calcula los laterales.
  \item Interpreta: asíntota vertical, horizontal, continuidad\ldots
\end{enumerate}

\begin{ojo}
\begin{itemize}
  \item \(\frac00\) no vale 0 ni 1: es una indeterminación.
  \item \(\frac k0\) exige mirar el signo a cada lado.
  \item \((fg)'\ne f'g'\) y \(\bigl(\frac fg\bigr)'\ne\frac{f'}{g'}\).
  \item Olvidar \(u'\) en la regla de la cadena.
  \item Aplicar L'Hôpital sin comprobar que hay \(\frac00\) o \(\frac\infty\infty\).
  \item Estudiar solo la continuidad en la unión de trozos y olvidar las derivadas laterales.
  \item Escribir la normal con la misma pendiente que la tangente.
\end{itemize}
\end{ojo}

\end{multicols}
\end{document}
